\documentclass[11pt,a4paper]{article}
\usepackage{fiz101notes}
\usepackage{needspace}
\sisetup{retain-explicit-plus=true}
\noteweek{2}{Kinematics: motion in one, two and three dimensions}

% Notation used only in this week (Young & Freedman style).
\newcommand{\vavx}{v_{\mathrm{av}\text{-}x}}   % average velocity, x component
\newcommand{\aavx}{a_{\mathrm{av}\text{-}x}}   % average acceleration, x component
\newcommand{\relx}[2]{#1_{#2\text{-}x}}        % relative quantity, x component: \relx{v}{P/A}

\begin{document}
\makenotetitle

%======================================================================
\setcounter{section}{2}
\section{Motion along a straight line}\label{sec:line}

Classical mechanics was born from the wish to describe and predict how bodies move, from falling
stones to the planets. The laws found by Galileo and Newton worked so well that a century later
the goal became very ambitious indeed: in 1814 Pierre-Simon Laplace argued that an intellect that
knew the forces of nature and the present positions and velocities of all the particles in the
universe could predict the future of the universe we live in and reconstruct its past, that is,
determine the history and the future of every particle in it. (We now know that this dream has
limits: quantum mechanics rules the world of atoms, Einstein's relativity takes over at speeds
close to the speed of light, and even within classical mechanics, in many systems tiny
uncertainties in the starting conditions grow so fast that long-term prediction is impossible in
practice. But for everyday objects moving at everyday speeds, the laws of classical mechanics are
extremely accurate.)

In Chapter~2 we introduced vectors, the basic language of classical mechanics. In this chapter we
take our last step before we start studying the dynamics of classical mechanics developed by
Galileo and Newton: we turn to \textbf{kinematics}, the description of motion, that is, of how the
positions of objects change with time. Kinematics asks \emph{how} things move; the question of
\emph{why} they move as they do, the forces, belongs to dynamics, which starts in Chapter~5. In
this chapter and the next we treat every moving object as a \textbf{particle}: an object so
small, compared with the distances involved, that its position at any time is a single point in
space. When we meet the center of mass in a later
chapter, we will see why this model works so well even for cars, cats and planets.

%----------------------------------------------------------------------
\subsection{Invariants}\label{sec:invariants}

\begin{figure}[ht]
\centering
\begin{minipage}[c]{0.58\textwidth}
\centering
\begin{tikzpicture}[x={(-0.62cm,-0.40cm)},y={(1.05cm,-0.10cm)},z={(0cm,1cm)},scale=1.5,
    >=Stealth,line cap=round]
  \coordinate (O) at (0,0,0);
  \coordinate (xp) at (0.866,-0.5,0);
  \coordinate (yp) at (0.441,0.765,-0.469);
  \coordinate (zp) at (0.235,0.407,0.883);
  % the xy plane
  \filldraw[fill=fzgray!12,draw=fzgray!70] (-0.9,-1.0,0) -- (1.25,-1.0,0) -- (1.25,1.25,0)
     -- (-0.9,1.25,0) -- cycle;
  % the x'y' plane
  \filldraw[fill=fzaux!10,draw=fzaux!60,opacity=0.85] ($(O)-0.9*(xp)-1.0*(yp)$)
     -- ($(O)+1.25*(xp)-1.0*(yp)$) -- ($(O)+1.25*(xp)+1.25*(yp)$) -- ($(O)-0.9*(xp)+1.25*(yp)$) -- cycle;
  % unprimed axes
  \draw[->,thick] (O) -- (1.55,0,0) node[below]{$x$};
  \draw[->,thick] (O) -- (0,1.6,0) node[right]{$y$};
  \draw[->,thick] (O) -- (0,0,1.45) node[above]{$z$};
  % primed axes
  \draw[->,thick,fzaux] (O) -- ($(O)+1.55*(xp)$) node[above left=-1pt]{$x'$};
  \draw[->,thick,fzaux] (O) -- ($(O)+1.6*(yp)$) node[right]{$y'$};
  \draw[->,thick,fzaux] (O) -- ($(O)+1.45*(zp)$) node[above]{$z'$};
  % the vector
  \draw[->,very thick,fzvec] (O) -- (0.35,0.95,1.05) node[above right=-1pt]{$\vec{A}$};
  % plane labels
  \node[circle,inner sep=0pt,pin={[fzgray,pin distance=6mm,pin edge={fzgray}]20:{$xy$ plane}}]
    at (-0.8,1.15,0) {};
  \node[circle,inner sep=0pt,pin={[fzaux,pin distance=6mm,pin edge={fzaux}]200:{$x'y'$ plane}}]
    at ($(O)+1.1*(xp)+1.1*(yp)$) {};
\end{tikzpicture}
\end{minipage}\hfill
\begin{minipage}[c]{0.38\textwidth}
\caption{Two coordinate systems, $xyz$ and $x'y'z'$, with a common origin; one is rotated with
respect to the other. The vector $\vec{A}$ is the same arrow in both, but its components are
different.}\label{fig:rotated}
\end{minipage}
\end{figure}

Before we start studying the kinematics of physical systems, it is worth stressing again a
property of vectors that we mentioned before: the laws of physics do not depend on how we choose
the coordinate axes. This also explains why we use vectors to describe motion. Let us look at the
length of a vector in two different coordinate systems. Let these two systems have the same
origin, but let one of them be rotated with respect to the other (Figure~\ref{fig:rotated}). In
these two coordinate systems, a vector $\vec{A}$ is written as
\begin{equation}\label{eq:A-two-systems}
  \vec{A} = A_x\ihat + A_y\jhat + A_z\khat , \qquad
  \vec{A} = A'_x\ihat' + A'_y\jhat' + A'_z\khat' .
\end{equation}
The components of $\vec{A}$ are different in the two coordinate systems: the same arrow is
described by different numbers. But $\vec{A}$ itself has not changed, it is the same arrow in
space, so its magnitude, the length of the arrow, must read the same in both coordinate systems:
\begin{equation}\label{eq:invariant-length}
  A_x^2 + A_y^2 + A_z^2 = A'^2_x + A'^2_y + A'^2_z .
\end{equation}
In other words, the magnitude of a vector is the same in every coordinate system; quantities of
this kind are called \emph{invariant} quantities. The components of a vector, on the other hand,
are not invariant: they change when we rotate the axes. Yet if two vectors are equal,
$\vec{A} = \vec{B}$, their components are equal in every coordinate system. So a law of physics
written as an equation between vectors holds in the same form whatever axes we choose, and this is
exactly why we write the laws of physics with vectors.

\begin{example}[Components change, the length does not]
Take $\vec{A} = 3\ihat + 4\jhat$ in the $xy$ plane, so $A = \sqrt{3^2+4^2} = 5$. Now rotate the
axes counterclockwise by an angle $\varphi$ about the $z$-axis. The new unit vectors are
$\ihat' = \cos\varphi\,\ihat + \sin\varphi\,\jhat$ and $\jhat' = -\sin\varphi\,\ihat +
\cos\varphi\,\jhat$, and the new components are the scalar products
\[
  A'_x = \vec{A}\cdot\ihat' = A_x\cos\varphi + A_y\sin\varphi , \qquad
  A'_y = \vec{A}\cdot\jhat' = -A_x\sin\varphi + A_y\cos\varphi .
\]
Squaring and adding, the cross terms $\pm 2A_xA_y\sin\varphi\cos\varphi$ cancel and
$\cos^2\varphi + \sin^2\varphi = 1$, so $A'^2_x + A'^2_y = A_x^2 + A_y^2$, as
\eqref{eq:invariant-length} says. For $\varphi = 30^\circ$: $A'_x = 3(0.866) + 4(0.500) = 4.60$
and $A'_y = -3(0.500) + 4(0.866) = 1.96$. The components have changed completely, but
$4.60^2 + 1.96^2 = 25.0 = 5^2$. (The scalar product $\vec{A}\cdot\vec{B} = AB\cos\theta$ of Week~1
is another invariant: it involves only the two lengths and the angle $\theta$ between the two
arrows.)
\end{example}

%----------------------------------------------------------------------
\subsection{Average and instantaneous velocity}\label{sec:velocity}

\begin{figure}[ht]
\centering
\begin{tikzpicture}[>=Stealth]
  \begin{scope}[rotate=15]
    \draw[thick] (0,0) -- (9,0);
    \foreach \p/\t/\x in {1.2/{t_1,\,t_5}/{x_1 = x_5}, 3.2/t_4/x_4, 5.1/t_2/x_2, 7.0/t_3/x_3}{
      \fill (\p,0) circle (1.6pt);
      \node at (\p,0.42) {$\t$};
      \node at (\p,-0.42) {$\x$};
    }
    \draw[->,thick,fzcomp] (0.2,0.85) -- (1.9,0.85) node[midway,above]{$x$};
  \end{scope}
\end{tikzpicture}
\caption{Motion in one dimension. The object moves along a straight line; the labels below the
line are its positions $x_1,\dots,x_5$ at the instants $t_1,\dots,t_5$ written above them. (It
passes some of these points more than once; for example, it also passes $x_4$ on the way out,
between $t_1$ and $t_2$.) The arrow marks the positive $x$ direction.}\label{fig:line}
\end{figure}

To study the kinematics of physical systems in one dimension, let us first look at the simple
example of Figure~\ref{fig:line}. An object moves along a straight line: it starts at position
$x_1$ at time $t_1$, passes $x_2$ at $t_2$ and $x_3$ at $t_3$, turns around a little beyond $x_3$,
passes $x_4$ at $t_4$, and at $t_5$ it is back at the position where it started. Motion of this
kind, along a single straight line, is called \textbf{one-dimensional motion}. Once we have chosen
an origin on the line and a positive direction (the arrow labeled $x$ in the figure), the position
of the object is given completely by one number, its coordinate $x$, which can be positive or
negative.

The \textbf{average velocity} of the object over any time interval is defined as
\begin{equation}\label{eq:vav}
  \vavx = \frac{x_2 - x_1}{t_2 - t_1} = \frac{\Delta x}{\Delta t} .
\end{equation}
The meaning of this definition is simple: the change of position between the time $t_1$ and the
time $t_2$, that is, the \emph{displacement} $\Delta x = x_2 - x_1$, divided by the time that has
passed, $\Delta t = t_2 - t_1$. Only the starting and the end points matter: what the object did
in between plays no role. (The subscript $x$ reminds us that this is the $x$ component of a
vector.) Returning to our example, we notice that the average velocity can be positive, negative
or zero. For example, between $t_1$ and $t_2$ the displacement $\Delta x$ is positive, so the
average velocity over this interval is positive; between $t_1$ and $t_5$ the displacement is zero,
so the average velocity is zero; and between $t_2$ and $t_4$ the displacement is negative, so the
average velocity is negative.

\begin{aside}
\textbf{Note: the choice of origin and of the $x$ direction.} From our example, you might conclude
intuitively that ``towards $x_3$'' is always the direction of positive average velocity. But what
decides in which intervals the average velocity is positive is the direction \emph{we} chose for
the $x$-axis. Had we chosen the $x$ direction the other way round, the displacement between $t_2$
and $t_4$ would be positive, and so would the average velocity over this interval. So keep in mind
that the direction of $x$ in the coordinate system you use to solve a problem determines the sign
of the average velocity. Remember this! Signs are very important in physics: a single minus sign
decides whether you owe me money or I owe you. Notice also that we have not chosen an origin at
all. The choice of origin plays no role in the average velocity: moving the origin adds the same
number to $x_1$ and to $x_2$, and this number cancels in the difference $x_2 - x_1$.
\end{aside}

\begin{figure}[ht]
\centering
\begin{tikzpicture}[x={(1.25cm,0cm)},y={(0cm,0.78cm)},>=Stealth]
  % axes
  \draw[->,very thick] (0,0) -- (6.0,0) node[right]{$t$};
  \draw[->,very thick] (0,0) -- (0,6.3) node[above]{$x$};
  % guide lines
  \foreach \t/\x in {1/1, 2/3.1, 3/4.3, 4/2, 5/1}{
    \draw[dashed,fzgray] (\t,0) -- (\t,\x);
  }
  \draw[dashed,fzgray] (0,1) -- (5,1);
  \draw[dashed,fzgray] (0,3.1) -- (2,3.1);
  \draw[dashed,fzgray] (0,4.3) -- (3,4.3);
  \draw[dashed,fzgray] (0,2) -- (4,2);
  \foreach \t/\l in {1/t_1, 2/t_2, 3/t_3, 4/t_4, 5/t_5}{
    \node[below] at (\t,0) {$\l$};
  }
  \node[left] at (0,1) {$x_1 = x_5$};
  \node[left] at (0,2) {$x_4$};
  \node[left] at (0,3.1) {$x_2$};
  \node[left] at (0,4.3) {$x_3$};
  % the x-t curve (piecewise cubic, see changes.md)
  \draw[thick,fzvec] (0,1) .. controls (0.3333,1) and (0.6667,1) .. (1,1)
    .. controls (1.2,1) and (1.4,1.3) .. (1.6,1.9)
    .. controls (1.7333,2.3) and (1.8667,2.7) .. (2,3.1)
    .. controls (2.3333,4.1) and (2.6667,4.1667) .. (3,4.3)
    .. controls (3.1333,4.3533) and (3.2667,5.3) .. (3.4,5.3)
    .. controls (3.6,5.3) and (3.8,3.7) .. (4,2)
    .. controls (4.04,1.66) and (4.08,1.53) .. (4.12,1.45)
    .. controls (4.2467,1.1967) and (4.3733,1.238) .. (4.5,1.2)
    .. controls (4.6667,1.15) and (4.8333,1.1167) .. (5,1)
    .. controls (5.1333,0.9067) and (5.2667,0.76) .. (5.4,0.6);
  % chord b-c with its triangle
  \draw[very thick,fzcomp] (2,3.1) -- (3,3.1) -- (3,4.3);
  \draw[very thick,fzcomp] (2,3.1) -- (3,4.3);
  \draw[fzcomp] (2.42,3.1) arc[start angle=0,end angle=36,x radius=0.42,y radius=0.54];
  \node[fzcomp] at (2.6,3.42) {$\alpha$};
  \node[below,fzcomp] at (2.5,3.1) {$\Delta t$};
  \node[right,fzcomp] at (3,3.7) {$\Delta x$};
  % chord d-e
  \draw[dashed,thick,fzcomp] (4,2) -- (5.2,2);
  \draw[dashed,thick,fzcomp] (4,2) -- (5,1);
  \draw[->,fzcomp] (4.55,2) arc[start angle=0,end angle=-44,x radius=0.55,y radius=0.7];
  \node[fzcomp] at (4.78,1.66) {$\beta$};
  % points
  \foreach \t/\x/\n/\pos in {1/1/a/above, 2/3.1/b/above left, 3/4.3/c/below right, 4/2/d/below left, 5/1/e/above right}{
    \fill (\t,\x) circle (1.5pt);
    \node[\pos] at (\t,\x) {$\n$};
  }
\end{tikzpicture}
\caption{Position--time graph for the motion of Figure~\ref{fig:line}. The average velocity
between two instants is the slope of the chord joining the two points of the curve: positive for
the rising chord $bc$ (angle $\alpha > 0$), negative for the falling chord $de$ (angle
$\beta < 0$).}\label{fig:xt}
\end{figure}

The position--time ($x$--$t$) graph in Figure~\ref{fig:xt} shows where the object of
Figure~\ref{fig:line} is at any time $t$. First, on this graph we have put the origin below $x_1$
(on the line of Figure~\ref{fig:line}, somewhere to the left of $x_1$). As we said, this choice makes no difference,
and any other point would do. Looking at this graph, if we take for example the change of position
between $t_2$ and $t_3$, the definition~\eqref{eq:vav} tells us that the average velocity is
positive. Geometrically, $\vavx$ is the \emph{slope} $\Delta x/\Delta t$ of the straight line, or
chord, that joins the points $b$ and $c$ of the curve. A positive average velocity means that the
chord rises, so the angle $\alpha$ it makes with the $t$ direction is positive. Similarly, between
$t_4$ and $t_5$ the average velocity is negative: the chord from $d$ to $e$ falls, and the
corresponding angle $\beta$, measured clockwise, is negative. On a given graph, the steeper the
chord, the larger the magnitude of the average velocity. (Strictly, $\tan\alpha$ is proportional
to the slope $\Delta x/\Delta t$, with a factor that depends on the scales we chose for the two
axes; the slope has units of \unit{m/s}, while an angle has none. So the angle is a picture of the
velocity, not the velocity itself.)

In everyday language we use the words \textbf{velocity} and \textbf{speed} as if they meant the
same thing. In physics, however, these two quantities are quite different, and to understand the
difference we need what we know about vectors and scalars. Over a time interval, the average speed
and the average velocity are defined as
\[
  \text{average speed} = \frac{\text{total distance traveled}}{\text{elapsed time}} , \qquad
  \text{average velocity} = \frac{\text{displacement}}{\text{elapsed time}} .
\]
There is a very important difference between distance traveled and displacement. The distance
traveled is a scalar and is never negative (if it is zero, you stood still the whole time). For
example, when you drive from one place to another, the distance traveled is the increase in the
reading of your car's odometer. The displacement is a vector, and like all vectors it has a
magnitude and a direction. The tail of the displacement vector is at the point where you started
and its tip at the point where you stopped; if you come back to your starting point, the
displacement is the zero vector $\vec{0}$, whatever path you took. In one dimension the
displacement vector is $\Delta x\,\ihat$, and the sign of its component $\Delta x$ gives its
direction. So for the interval from $t_1$ to $t_5$ in our example, the average velocity is zero,
but the average speed is not. In general, the average speed is \emph{not} the magnitude of the
average velocity.

\begin{figure}[ht]
\centering
\begin{tikzpicture}[>=Stealth]
  % axes (t2 at x=1, t3 at x=6; x2 at y=0.78, x3 at y=3.9)
  \draw[->,very thick] (0,0) -- (7,0) node[right]{$t$};
  \draw[->,very thick] (0,0) -- (0,4.6) node[above]{$x$};
  \draw[dashed,fzgray] (0,0.78) -- (1,0.78);
  \draw[dashed,fzgray] (0,3.9) -- (6,3.9);
  \node[left] at (0,0.78) {$x_2$};
  \node[left] at (0,3.9) {$x_3$};
  \draw[dashed,fzgray] (1,0) -- (1,0.78);
  \draw[dashed,fzgray] (2.25,0) -- (2.25,2.316);
  \draw[dashed,fzgray] (3.5,0) -- (3.5,3.185);
  \draw[dashed,fzgray] (4.75,0) -- (4.75,3.632);
  \draw[dashed,fzgray] (6,0) -- (6,3.9);
  \node[below] at (1,0) {$t_2$};
  \node[below] at (6,0) {$t_3$};
  % chords from b
  \draw[thick,fzcomp] (1,0.78) -- (6,0.78) -- (6,3.9);
  \draw[very thick,fzcomp] (1,0.78) -- (6,3.9);
  \draw[dashed,fzcomp] (1,0.78) -- (4.75,3.632);
  \draw[dashed,fzcomp] (1,0.78) -- (3.5,3.185);
  \draw[dashed,fzcomp] (1,0.78) -- (2.25,2.316);
  % tangent at b
  \draw[thick,fzaux] (0.55,0.078) -- (2.6,3.276) node[above,font=\small]{tangent at $b$};
  % the curve
  \draw[thick,fzvec] (1,0.78) .. controls (2.667,3.38) and (4.333,3.553) .. (6,3.9);
  \draw[fzcomp] (1.75,0.78) arc[start angle=0,end angle=32,radius=0.75];
  \node[fzcomp] at (2.05,1.02) {$\alpha$};
  \draw[->,thick] (5.4,0.35) -- (4.2,0.35);
  \fill (1,0.78) circle (1.5pt) node[above left]{$b$};
  \fill (6,3.9) circle (1.5pt) node[above right]{$c$};
\end{tikzpicture}
\caption{Instantaneous velocity. As $t_3$ moves towards $t_2$ (arrow), the chord from $b$ to the
curve (dashed lines) turns and approaches the tangent to the curve at $b$. The slope of this
tangent is the instantaneous velocity at $t_2$.}\label{fig:tangent}
\end{figure}

Now let us look at Figure~\ref{fig:tangent} to study how the position changes between $t_2$ and
$t_3$ in more detail. As the figure shows, if we move $t_3$ closer and closer to $t_2$, the chord
from $b$ turns and, for this curve, the angle $\alpha$ grows. In the limit as $t_3$ approaches
$t_2$ ($\Delta t \to 0$), the chord approaches the \emph{tangent} to the $x$--$t$ curve at $b$, and
its slope approaches the slope of this tangent. (We cannot simply set $t_3 = t_2$: then
$\Delta x = 0$ and $\Delta t = 0$, and $0/0$ has no meaning; only the limit does.) This gives us the definition of the \textbf{instantaneous velocity} at a
time $t$:
\begin{equation}\label{eq:vinst-limit}
  v_x = \lim_{\Delta t\to 0}\frac{x(t+\Delta t) - x(t)}{\Delta t} = \dv{x}{t} .
\end{equation}
This formula, which you should keep in mind for this course and for the rest of your life, is
worth writing once more:
\begin{equation}\label{eq:vinst}
  \boxed{\,v_x = \dv{x}{t}\,} .
\end{equation}
In words: the instantaneous velocity is the rate of change of position with time, and on an
$x$--$t$ graph it is the slope of the tangent to the curve at that instant. From now on,
``velocity'' without an adjective means the instantaneous velocity.

This quantity can be negative, zero or positive, according to whether the tangent falls, is
horizontal or rises (that is, whether its angle with the $t$ direction is negative, zero or
positive). So, as we can see in Figure~\ref{fig:xt}, the velocity of the object is zero between
$0$ and $t_1$, where the object is at rest, and at the turning point between $t_3$ and $t_4$ (the
top of the curve); it is positive from $t_1$ until this turning point, and negative from the
turning point until $t_5$ and after. The \textbf{instantaneous speed} is the magnitude of the
instantaneous velocity. Like the average speed, it does not depend on our choice of direction, and
it is never negative. For example, if the velocity of the object is \qty{30}{m/s} at $t_2$ and
\qty{-100}{m/s} at $t_4$, then $v_x$ at $t_2$ is the larger of the two as a signed number
($+30 > -100$), but the speed at $t_2$ (\qty{30}{m/s}) is smaller than the speed at $t_4$
(\qty{100}{m/s}). (Velocities themselves are vectors, and vectors cannot be ordered as larger or
smaller; only their components and their magnitudes can.)

%----------------------------------------------------------------------
\subsection{Average and instantaneous acceleration}\label{sec:acceleration}

Just as velocity gives the change of position per unit time, \textbf{acceleration} gives the change
of velocity per unit time. The average acceleration over a time interval is defined as
\begin{equation}\label{eq:aav}
  \aavx = \frac{v_{2x} - v_{1x}}{t_2 - t_1} = \frac{\Delta v_x}{\Delta t} .
\end{equation}
Like the average velocity, the average acceleration can be negative, zero or positive. For
example, in the motion of Figure~\ref{fig:xt} the velocity is zero between $0$ and $t_1$, and at
$t_1$ the object starts to move: the slope of the curve starts to grow. Between $t_1$ and $t_2$
the velocity increases from zero to a positive value, so the average acceleration over this
interval is positive. But if we look at the average acceleration over the interval
from $t_1$ to $t_5$, it is negative, because the velocity is zero at $t_1$ and negative at $t_5$.
As these examples show, both the average velocity and the average acceleration depend on how we
choose the direction of increasing $x$: if we reverse it, the signs of both change. (On a graph of
$v_x$ against $t$, the average acceleration is the slope of the chord between the two points, just
as the average velocity is the slope of a chord on the $x$--$t$ graph.)

\begin{figure}[ht]
\centering
\begin{minipage}[c]{0.46\textwidth}
\centering
\begin{tikzpicture}[>=Stealth]
  \fill[pattern=north east lines,pattern color=fzgray] (-1.5,-0.3) rectangle (1.5,0);
  \draw[thick] (-1.5,0) -- (1.5,0);
  \draw[dashed,fzgray] (-1.1,2.2) -- (1.1,2.2);
  \draw[->,very thick,fzvec] (-0.6,2.2) -- (-0.6,0.05);
  \draw[->,very thick,fzvec] (0.6,0.05) -- (0.6,2.2);
  \node[above,font=\small] at (-0.6,2.2) {before};
  \node[above,font=\small] at (0.6,2.2) {after};
  \node[left,fzvec] at (-0.6,1.1) {$v_x = \qty{-5}{m/s}$};
  \node[right,fzvec] at (0.6,1.1) {$v_x = \qty{+5}{m/s}$};
  \draw[->,thick,fzcomp] (-3.4,0.5) -- (-3.4,1.7) node[above]{$+x$};
\end{tikzpicture}
\end{minipage}\hfill
\begin{minipage}[c]{0.44\textwidth}
\caption{The velocity of a tennis ball just before and just after it bounces off the ground;
upward is the positive $x$ direction.}\label{fig:bounce}
\end{minipage}
\end{figure}

Now consider a tennis ball that hits the ground with velocity \qty{-5}{m/s} and bounces back up with
velocity \qty{+5}{m/s}, as in Figure~\ref{fig:bounce}; here we take upward as the positive $x$
direction. To know the average acceleration we need to know over how much time this change in
velocity happens. If we assume that the ball is in contact with the ground for
$\Delta t = \qty{e-2}{s}$, the average acceleration is
\begin{equation}\label{eq:tennis}
  \aavx = \frac{\qty{5}{m/s} - (\qty{-5}{m/s})}{\qty{e-2}{s}} = \qty{1000}{m/s^2} ,
\end{equation}
about a hundred times the acceleration of a falling body. If we reverse the $x$ direction (and
take downward as positive), the sign of the average acceleration reverses too, and
$\aavx = \qty{-1000}{m/s^2}$. You can also try this experiment with a tomato or an egg, but then
the objects thrown at the ground with \qty{-5}{m/s} will not bounce back: their velocity just drops
to zero. The change in velocity is then \qty{5}{m/s}, not \qty{10}{m/s}. If the squashing lasts,
say, \qty{0.02}{s}, the average acceleration is $\aavx = (\qty{5}{m/s})/(\qty{0.02}{s}) =
\qty{250}{m/s^2}$. (Is \qty{0.02}{s} reasonable? While it stops, the tomato moves on average at
about \qty{2.5}{m/s}, so in \qty{0.02}{s} it squashes by about \qty{5}{cm}, which is about right.
A stop lasting \qty{0.25}{s} would need some \qty{60}{cm}.) Here again the direction of $x$ decides the sign of the average acceleration.
But remember that the egg or the tomato does not care which direction you chose as the $x$-axis.
In every one of these cases the acceleration \emph{vector} during the impact points upward (the
ground pushes the ball up); only the sign of its component depends on our axis. Physics does not
depend on such choices; the choices you make are only tools that help you describe physical
phenomena.

Just as we went from the average velocity to the velocity at an instant, we can go from the
average acceleration to the \textbf{instantaneous acceleration}:
\begin{equation}\label{eq:ainst-limit}
  a_x = \lim_{\Delta t\to 0}\frac{v_x(t+\Delta t) - v_x(t)}{\Delta t} = \dv{v_x}{t}
      = \frac{\dd^2 x}{\dd t^2} .
\end{equation}
The last form says that acceleration is the \emph{second derivative} of position: we differentiate
$x(t)$ once to get $v_x$, and once more to get $a_x$. This second formula, which should also stay
with you for the rest of this course and of your life, is worth writing again:
\begin{equation}\label{eq:ainst}
  \boxed{\,a_x = \dv{v_x}{t} = \frac{\dd^2 x}{\dd t^2}\,} .
\end{equation}
Just as $v_x$ is the slope of the tangent to the $x$--$t$ curve, $a_x$ is the slope of the tangent
to the $v_x$--$t$ curve.

Now let us return to our example in Figure~\ref{fig:xt} to find the instantaneous acceleration at
various times. For the velocity we had a good intuitive tool: whether the tangent rises or falls
told us the sign of the velocity. For the acceleration we must look not at the angle of the
tangent itself but at how this angle \emph{changes}, since $a_x$ is the rate at which the slope of
the $x$--$t$ curve changes. Where the slope increases, the curve bends upward (it is concave up)
and $a_x > 0$; where the slope decreases, the curve bends downward (concave down) and $a_x < 0$;
where the graph is a straight line, $a_x = 0$. Between $0$ and $t_1$ the velocity, that is the
slope, does not change, so the acceleration is zero in this interval. Between $t_1$ and $t_2$ the
velocity first increases and then stays constant (the curve becomes a straight line just before
$t_2$), so the object first has a positive acceleration and then zero acceleration. Between $t_3$
and $t_4$ the acceleration is first positive (just after $c$ the curve bends upward) and then
negative (the curve bends downward over the top and on the way down to $d$).

\begin{figure}[ht]
\centering
\begin{tikzpicture}
\begin{axis}[name=xt,width=0.48\textwidth,height=0.42\textwidth,axis lines=left,
    xmin=0,xmax=6.6,ymin=-2,ymax=10.5,xlabel={$t$ (s)},ylabel={$x$ (m)},
    xtick={0,1,2,3,4,5,6},ytick={-1,0,2,4,6,8,10},tick label style={font=\footnotesize},
    label style={font=\small},title={(a)},title style={font=\small}]
  \addplot[fzgray,domain=0:6.6,samples=2] {0};
  \addplot[thick,fzvec,domain=0:6.3,samples=80] {8-6*x+x^2};
  \addplot[dashed,fzgray] coordinates {(0,-1) (3,-1) (3,-2)};
  \addplot[only marks,mark=*,mark size=1.4pt] coordinates {(0,8) (2,0) (3,-1) (4,0)};
\end{axis}
\begin{axis}[at={(xt.east)},anchor=west,xshift=0.6cm,width=0.48\textwidth,height=0.42\textwidth,
    axis lines=left,xmin=0,xmax=6.6,ymin=-7,ymax=7.5,xlabel={$t$ (s)},ylabel={$v_x$ (m/s)},
    xtick={0,1,2,3,4,5,6},ytick={-6,-4,-2,0,2,4,6},tick label style={font=\footnotesize},
    label style={font=\small},title={(b)},title style={font=\small}]
  \addplot[fzgray,domain=0:6.6,samples=2] {0};
  \addplot[thick,fzvec,domain=0:6.3,samples=2] {-6+2*x};
  \addplot[only marks,mark=*,mark size=1.4pt] coordinates {(0,-6) (3,0) (4,2)};
\end{axis}
\end{tikzpicture}
\caption{The object with position $x = (8 - 6t + t^2)\,\unit{m}$. (a) Position--time graph: the slope of
the curve is the velocity. (b) Velocity--time graph, $v_x = (-6 + 2t)\,\unit{m/s}$: its constant slope
is the acceleration $a_x = \qty{2}{m/s^2}$.}\label{fig:parabola}
\end{figure}

As an example, consider an object whose position is $x = (8 - 6t + t^2)\,\unit{m}$, with $t$ in seconds.
(More carefully, $x = \qty{8}{m} - (\qty{6}{m/s})\,t + (\qty{1}{m/s^2})\,t^2$, so that every term
is in meters.) Recall that if $x = t^n$, then $\dd x/\dd t = nt^{n-1}$, and that the derivative of
a constant is zero. Then the velocity of the object at any time $t$ is
\begin{equation}\label{eq:ex-v}
  v_x = (-6 + 2t)\,\unit{m/s} ,
\end{equation}
and its acceleration at any time $t$ is
\begin{equation}\label{eq:ex-a}
  a_x = \qty{2}{m/s^2} .
\end{equation}
So in this case the acceleration does not change with time: it is constant. With these results, at
$t = 0$ we have
\begin{equation}\label{eq:ex-t0}
  t = 0 , \qquad x = \qty{+8}{m} , \qquad v_x = \qty{-6}{m/s} , \qquad a_x = \qty{2}{m/s^2} .
\end{equation}
The object passes through the origin $x = 0$ twice, at the two roots of
$8 - 6t + t^2 = (t-2)(t-4) = 0$: once at $t = \qty{2}{s}$ and once at $t = \qty{4}{s}$. To find
where the object stops and turns around, we need to find when its velocity is zero. From
$-6 + 2t = 0$, the velocity is zero at $t = \qty{3}{s}$. Substituting $t = \qty{3}{s}$ into the
position equation, the turning point is at $x = 8 - 18 + 9 = \qty{-1}{m}$.

The motion is now clear. As Figure~\ref{fig:parabola} shows, the object starts at $x = \qty{8}{m}$
with velocity $v_x = \qty{-6}{m/s}$ and acceleration $a_x = \qty{2}{m/s^2}$, then passes $x = 0$,
reaches $x = \qty{-1}{m}$, where its velocity is zero, turns around, and passes $x = 0$ again at
$t = \qty{4}{s}$. These results agree with the graph. For example, between $t = 0$ and
\qty{3}{s} the slope is negative and, as expected, so is the velocity. But since the acceleration
is positive, we expect the velocity to become zero at some point and then positive. In line with
this expectation, at $t = \qty{3}{s}$ the slope is zero and so is the velocity. After
$t = \qty{3}{s}$ the slope, and so the velocity, is positive. For example, at $t = \qty{4}{s}$ we
expect a positive velocity, and the velocity equation~\eqref{eq:ex-v} tells us that
$v_x = \qty{+2}{m/s}$ at this instant.

Two points are worth stressing. First, at $t = \qty{3}{s}$ the velocity is zero but the
acceleration is not: it is \qty{2}{m/s^2} at all times. An object can be at rest for an instant and
still be accelerating; that is precisely why it does not stay at rest. Second, the parabola in
Figure~\ref{fig:parabola}(a) is a graph of $x$ against $t$, not the path of the object: the object
moves back and forth along the straight $x$-axis. To see this example in a simulation, I recommend
\href{https://phet.colorado.edu/en/simulations/moving-man}{\emph{The Moving Man}},\footnote{\url{https://phet.colorado.edu/en/simulations/moving-man}}
a PhET simulation developed at the University of Colorado.

\begin{keyresult}
\textbf{Speeding up or slowing down?} The sign of $a_x$ alone does not tell. An object speeds up
when $v_x$ and $a_x$ have the \emph{same} sign and slows down when they have \emph{opposite} signs.
In the example above, for $0 < t < \qty{3}{s}$ we have $v_x < 0$ and $a_x > 0$: the object slows
down although its acceleration is positive. For $t > \qty{3}{s}$ both are positive, and it speeds
up.
\end{keyresult}

For an object moving in one dimension with \emph{constant} acceleration, the most general position,
velocity and acceleration are, with $c_1$, $c_2$ and $c_3$ any constants,
\begin{equation}\label{eq:general-const-a}
\begin{aligned}
  x   &= c_1 + c_2 t + c_3 t^2 , \\
  v_x &= c_2 + 2c_3 t , \\
  a_x &= 2c_3 .
\end{aligned}
\end{equation}
Why is this the most general form? The velocity is a function whose derivative is the constant
$a_x$, so it must be a straight line in $t$; the position is a function whose derivative is this
straight line, so it must be a quadratic in $t$. (In the language of calculus: we integrate $a_x$
once to get $v_x$ and once more to get $x$, and each integration brings in one constant.) If we
write $x_0$ and $v_{0x}$ for the position and the velocity at $t = 0$, putting $t = 0$ into
\eqref{eq:general-const-a} gives $c_1 = x_0$ and $c_2 = v_{0x}$, and the last line gives
$c_3 = \tfrac12 a_x$. Therefore, for motion with constant acceleration the most general position
and velocity equations are
\begin{equation}\label{eq:const-a}
  x = x_0 + v_{0x}t + \tfrac12 a_x t^2 , \qquad v_x = v_{0x} + a_x t .
\end{equation}
By playing with these equations we can obtain various results that are useful for straight-line
motion with constant acceleration. For example, using $a_x t = v_x - v_{0x}$ from the second
equation, we can write the first one as
\begin{equation}\label{eq:const-a-avg}
  x - x_0 = v_{0x}t + \tfrac12 a_x t^2 = v_{0x}t + \tfrac12 (v_x - v_{0x})t
          = \tfrac12 (v_{0x} + v_x)\,t .
\end{equation}
This says that when the acceleration is constant, the average velocity $(x - x_0)/t$ is simply the
mean of the initial and final velocities; this is not true for a general motion. Or we can square
the velocity equation and use $x - x_0 = v_{0x}t + \tfrac12 a_x t^2$ from the first one:
\begin{equation}\label{eq:const-a-v2}
\begin{split}
  v_x^2 &= (v_{0x} + a_x t)^2 = v_{0x}^2 + 2v_{0x}a_x t + a_x^2 t^2 \\
        &= v_{0x}^2 + 2a_x\bigl(v_{0x}t + \tfrac12 a_x t^2\bigr) = v_{0x}^2 + 2a_x (x - x_0) .
\end{split}
\end{equation}
This last equation does not contain the time: it is the one to use when the time is neither given
nor asked for.

\begin{keyresult}
\textbf{Straight-line motion with constant acceleration} (valid \emph{only} if $a_x$ is constant):
\[
\begin{aligned}
  v_x &= v_{0x} + a_x t , &\qquad x &= x_0 + v_{0x}t + \tfrac12 a_x t^2 , \\
  v_x^2 &= v_{0x}^2 + 2a_x(x - x_0) , &\qquad x - x_0 &= \tfrac12(v_{0x} + v_x)\,t .
\end{aligned}
\]
If the acceleration changes with time, these equations are wrong, and we must go back to the
definitions \eqref{eq:vinst} and \eqref{eq:ainst}.
\end{keyresult}

\needspace{8\baselineskip}
\begin{example}[Stopping distance]
A car moving at $v_{0x} = \qty{20}{m/s}$ (\qty{72}{km/h}) brakes with constant acceleration
$a_x = \qty{-5.0}{m/s^2}$. How far does it go before it stops? The time is not asked, so we use
\eqref{eq:const-a-v2} with $v_x = 0$:
\[
  0 = v_{0x}^2 + 2a_x(x - x_0) \quad\Rightarrow\quad
  x - x_0 = -\frac{v_{0x}^2}{2a_x} = \frac{(\qty{20}{m/s})^2}{2(\qty{5.0}{m/s^2})} = \qty{40}{m} .
\]
At twice the speed, \qty{40}{m/s}, the stopping distance is four times as long, \qty{160}{m}: the
stopping distance grows as the \emph{square} of the speed. (Real stopping distances are longer
still, because the driver needs some reaction time before the brakes are applied.)
\end{example}

%----------------------------------------------------------------------
\subsection{Free fall}\label{sec:freefall}

Now that we have learned the basics of straight-line motion, we can come back to bodies falling
freely under gravity, which we studied by dimensional analysis in Section~1.1.1 (Week~1).

\begin{figure}[ht]
\centering
\begin{minipage}[c]{0.40\textwidth}
\centering
\begin{tikzpicture}[>=Stealth]
  \fill[pattern=north east lines,pattern color=fzgray] (-1.4,-0.3) rectangle (1.4,0);
  \draw[thick] (-1.4,0) -- (1.4,0);
  \draw[dashed,fzgray] (-1.2,3) -- (1.2,3) node[right,black]{$x = 0$};
  \fill[fzvec] (0,3.18) circle (0.18);
  \draw[->,fzgray] (0,2.95) -- (0,0.05);
  \draw[<->] (-1.0,3) -- (-1.0,0) node[midway,left]{$h$};
  \draw[->,very thick,fzaux] (0.45,2.2) -- (0.45,1.3) node[right]{$\vec{a}$};
  \draw[->,thick,fzcomp] (2.2,2.4) -- (2.2,1.2) node[below]{$+x$};
\end{tikzpicture}
\end{minipage}\hfill
\begin{minipage}[c]{0.50\textwidth}
\caption{Free fall from rest from a height $h$. Here the positive $x$ direction is downward, so
$a_x = +g$.}\label{fig:freefall}
\end{minipage}
\end{figure}

The best-known example of motion with constant acceleration is the motion of bodies falling
freely under the Earth's gravity. In Section~1.1.1 dimensional analysis showed that the time of
fall cannot depend on the mass of the body, and we saw that this rests on an experimental fact
first established by Galileo: in a vacuum, where there is no air resistance, all bodies at the
same place fall with the same acceleration, whatever their mass, chemical composition or shape. (In
air, the shape and the size matter, through air resistance; that is why a feather falls slowly.)
We call \textbf{free fall} any motion under the influence of gravity alone, with air resistance
neglected: a body dropped from rest, a body thrown downward and a body thrown upward are all in
free fall, on the way up as well as on the way down. The magnitude of the acceleration due to
gravity, a positive number, is denoted by $g$ and, to two significant figures, is
\begin{equation}\label{eq:g}
  g = \qty{9.8}{m/s^2} .
\end{equation}
The acceleration vector of a freely falling body always points vertically downward (very nearly
towards the center of the Earth); whether its component is $+g$ or $-g$ depends on which direction we call positive. (The
value of $g$ varies slightly from place to place, from about \qty{9.78}{m/s^2} at the equator to
about \qty{9.83}{m/s^2} at the poles, and it decreases with altitude. Treating it as a constant is a
good model as long as the heights involved are small compared with the radius of the Earth,
\qty{6370}{km}, and air resistance is negligible.)

\begin{aside}
To see that a feather and a bowling ball really fall together when there is no air, watch
\href{https://www.youtube.com/watch?v=E43-CfukEgs}{feathers and a bowling ball dropped in NASA's
vacuum chamber}: \url{https://www.youtube.com/watch?v=E43-CfukEgs}.
\end{aside}

Consider a body that is at position $x_0 = 0$ with velocity $v_{0x} = 0$ at time $t = 0$, as in
Figure~\ref{fig:freefall}. If we choose the direction towards the ground (downward) as positive,
then $a_x = +g$, and the first equation of \eqref{eq:const-a} gives $x = \frac12 gt^2$. The body
reaches the ground when $x = h$, so the time it takes to fall is
\begin{equation}\label{eq:fall-time}
  t_h = \sqrt{\frac{2h}{g}} .
\end{equation}
So the constant $C$ in equation~(1.9) that dimensional analysis could not determine in
Section~1.1.1 is $C = \sqrt{2}$.

\begin{example}[Dropping a ball, and throwing it up]
(a) A ball is dropped from rest from a height of \qty{2.0}{m}, and then from \qty{8.0}{m}. From
\eqref{eq:fall-time}, $t_h = \sqrt{2(\qty{2.0}{m})/(\qty{9.8}{m/s^2})} = \qty{0.64}{s}$ and
$\sqrt{2(\qty{8.0}{m})/(\qty{9.8}{m/s^2})} = \qty{1.28}{s}$: four times the height, but only twice
the time, as Week~1 promised. From \eqref{eq:const-a-v2} with $v_{0x} = 0$ and $a_x = g$, the ball
reaches the ground with speed $v = \sqrt{2gh}$, that is \qty{6.3}{m/s} and \qty{12.5}{m/s}.

(b) Now throw the ball straight up at \qty{9.8}{m/s}. This time let us call the vertical axis $y$
and take \emph{upward} as positive, so $v_{0y} = \qty{+9.8}{m/s}$ and $a_y = -g =
\qty{-9.8}{m/s^2}$. The velocity $v_y = v_{0y} - gt$ is zero at $t = \qty{1.0}{s}$, when the ball
is at its highest point, $v_{0y}^2/(2g) = \qty{4.9}{m}$ above the hand (from
\eqref{eq:const-a-v2} with $v_y = 0$). At that instant the velocity is zero but the acceleration
is still \qty{-9.8}{m/s^2}: the ball is in free fall all the time, on its way up, at the top and on
its way down. (If its acceleration were zero at the top, it would simply stay there.) By symmetry
it is back in the hand at $t = \qty{2.0}{s}$, moving downward at \qty{9.8}{m/s}.
\end{example}

%======================================================================
\setcounter{section}{3}
\section{Motion in two and three dimensions}\label{sec:3d}

The aim of this chapter is to combine vectors, which we studied in Chapter~2, with motion in one
dimension, which we studied in Chapter~3, to describe motion in two and three dimensions. For this
we first need to be able to give the position of the object whose motion we want to study. We
describe the position of an object at a point $P = (x, y, z)$ of three-dimensional space by
its \textbf{position vector}, the vector from the origin to that point:
\begin{equation}\label{eq:r}
  \vec{r} = x\ihat + y\jhat + z\khat .
\end{equation}
As the object moves, its coordinates change with time: they are functions $x(t)$, $y(t)$ and
$z(t)$.

To study the kinematics of this object, we can follow the procedure of the previous chapter. First,
if the object has moved from the point $P_1 = (x_1, y_1, z_1)$ to the point $P_2 = (x_2, y_2, z_2)$
in a time interval $\Delta t$, its average velocity is
\begin{equation}\label{eq:vav3}
  \avg{\vec{v}} = \frac{\Delta\vec{r}}{\Delta t} ,
\end{equation}
where
\begin{equation}\label{eq:dr}
  \Delta\vec{r} = \vec{r}_2 - \vec{r}_1 = (x_2 - x_1)\ihat + (y_2 - y_1)\jhat + (z_2 - z_1)\khat
\end{equation}
is the displacement vector. As in one dimension, the average velocity depends only on the
displacement, not on the path between $P_1$ and $P_2$. The instantaneous velocity is again defined
by the limit $\Delta t\to 0$:
\begin{equation}\label{eq:v3}
  \vec{v} = \lim_{\Delta t\to 0}\frac{\Delta\vec{r}}{\Delta t} = \dv{\vec{r}}{t} .
\end{equation}
As $\Delta t \to 0$, the displacement $\Delta\vec{r}$ points more and more along the path, so the
instantaneous velocity is always \emph{tangent to the path}; its magnitude $v = |\vec{v}|$ is the
speed. Similarly, we define the average acceleration for motion in three dimensions as
\begin{equation}\label{eq:aav3}
  \avg{\vec{a}} = \frac{\Delta\vec{v}}{\Delta t} ,
\end{equation}
and the instantaneous acceleration as
\begin{equation}\label{eq:a3}
  \vec{a} = \lim_{\Delta t\to 0}\frac{\Delta\vec{v}}{\Delta t} = \dv{\vec{v}}{t}
          = \frac{\dd^2\vec{r}}{\dd t^2} .
\end{equation}
Here $\Delta\vec{v} = \vec{v}_2 - \vec{v}_1$ is a vector difference: the velocity changes when its
magnitude changes, when its direction changes, or both. So an object moving along a curved path
always has an acceleration, even if its speed is constant. We will meet the most important example,
uniform circular motion, in Section~\ref{sec:circular}.

If we use the components of $\vec{r}$ as in \eqref{eq:r}, we can write the velocity and
acceleration vectors in components as follows:
\begin{equation}\label{eq:components3}
  \vec{r} = x\ihat + y\jhat + z\khat , \qquad
  \vec{v} = \dv{x}{t}\ihat + \dv{y}{t}\jhat + \dv{z}{t}\khat , \qquad
  \vec{a} = \frac{\dd^2x}{\dd t^2}\ihat + \frac{\dd^2y}{\dd t^2}\jhat + \frac{\dd^2z}{\dd t^2}\khat .
\end{equation}
Here we used the fact that the unit vectors $\ihat$, $\jhat$, $\khat$ point along fixed axes and do
not change with time, so only the components need to be differentiated. In other words,
$v_x = \dd x/\dd t$, $v_y = \dd y/\dd t$, $v_z = \dd z/\dd t$, and $a_x = \dd v_x/\dd t$, and so on.

These equations show us a very simple but very important result. Each component of the velocity
and of the acceleration involves only the coordinate along the same axis: $v_x$ and $a_x$ come from
$x(t)$ alone, $v_y$ and $a_y$ from $y(t)$ alone, and so on. So, as long as each component of the
acceleration is known as a function of time alone (for example, when the acceleration is
constant), a three-dimensional problem reduces to \emph{three independent one-dimensional
problems}, one along each axis; the only thing they share is the time $t$. (In later chapters the
components can be tied together through the forces: the air resistance on a fast ball, for
example, grows roughly as the square of its total speed, so its horizontal deceleration depends on
how fast it is also moving vertically. For the constant accelerations of this chapter the
components are completely independent.) If the acceleration is constant, we can use the
one-dimensional equations \eqref{eq:const-a} for each axis, and the position and velocity along each
axis are
\begin{equation}\label{eq:const-a3}
\begin{aligned}
  x &= x_0 + v_{0x}t + \tfrac12 a_x t^2 , &\qquad v_x &= v_{0x} + a_x t , \\
  y &= y_0 + v_{0y}t + \tfrac12 a_y t^2 , &\qquad v_y &= v_{0y} + a_y t , \\
  z &= z_0 + v_{0z}t + \tfrac12 a_z t^2 , &\qquad v_z &= v_{0z} + a_z t .
\end{aligned}
\end{equation}
Here $a_x$, $a_y$ and $a_z$ are the (constant) components of the acceleration along $x$, $y$ and
$z$, and $x_0$, $v_{0x}$ and so on are the components of the position and velocity at $t = 0$.

%----------------------------------------------------------------------
\subsection{Projectile motion}\label{sec:projectile}

A good example of how motion in more than one dimension reduces to a sum of one-dimensional motions
is projectile motion.

\begin{figure}[ht]
\centering
\begin{tikzpicture}[x={(1.12cm,0cm)},y={(0cm,1.12cm)},>=Stealth]
  % axes
  \draw[->,thick] (0,0) -- (10.9,0) node[right]{$x$};
  \draw[->,thick] (0,0) -- (0,4.3) node[above]{$y$};
  \node[below left] at (0,0) {$O$};
  % trajectory
  \draw[dashed,fzgray,thick,domain=0:9,samples=60] plot (\x,{1.5*\x-0.16667*\x*\x});
  % launch: v0 = (1.28,1.92)
  \draw[dashed,fzgray] (1.28,0) -- (1.28,1.92) -- (0,1.92);
  \draw[->,thick,dashed,fzcomp] (0,0) -- (1.28,0) node[below]{$v_{0x}$};
  \draw[->,thick,dashed,fzcomp] (0,0) -- (0,1.92) node[left]{$v_{0y}$};
  \draw[->,very thick,fzvec] (0,0) -- (1.28,1.92) node[above left]{$\vec{v}_0$};
  \draw (0.75,0) arc[start angle=0,end angle=56.3,radius=0.75];
  \node at (0.95,0.5) {$\theta$};
  % rising point (2.25, 2.531), v = (1.28, 0.96)
  \draw[->,thick,dashed,fzcomp] (2.25,2.531) -- (3.53,2.531) node[right]{$v_x$};
  \draw[->,thick,dashed,fzcomp] (2.25,2.531) -- (2.25,3.491) node[above]{$v_y$};
  \draw[->,very thick,fzvec] (2.25,2.531) -- (3.53,3.491) node[above]{$\vec{v}$};
  % top (4.5, 3.375), v = (1.28, 0)
  \draw[->,very thick,fzvec] (4.5,3.375) -- (5.78,3.375) node[above]{$\vec{v} = v_x\ihat$};
  % falling point (6.75, 2.531), v = (1.28,-0.96)
  \draw[->,thick,dashed,fzcomp] (6.75,2.531) -- (8.03,2.531) node[right]{$v_x$};
  \draw[->,thick,dashed,fzcomp] (6.75,2.531) -- (6.75,1.571) node[below]{$v_y$};
  \draw[->,very thick,fzvec] (6.75,2.531) -- (8.03,1.571) node[right]{$\vec{v}$};
  % landing (9,0), v = (1.28,-1.92)
  \draw[->,thick,dashed,fzcomp] (9,0) -- (10.28,0) node[above]{$v_x$};
  \draw[->,thick,dashed,fzcomp] (9,0) -- (9,-1.92) node[left]{$v_y$};
  \draw[->,very thick,fzvec] (9,0) -- (10.28,-1.92) node[right]{$\vec{v}$};
  \foreach \p in {(0,0),(2.25,2.531),(4.5,3.375),(6.75,2.531),(9,0)} \fill \p circle (1.5pt);
  % acceleration
  \draw[->,very thick,fzaux] (4.5,1.5) -- (4.5,0.6) node[right]{$\vec{a} = -g\jhat$};
\end{tikzpicture}
\caption{Projectile motion. The horizontal velocity component $v_x$ is the same everywhere; the
vertical component $v_y$ decreases at the steady rate $g$, is zero at the highest point and
negative on the way down. The acceleration is $-g\jhat$ at every point of the path.}
\label{fig:projectile}
\end{figure}

A \textbf{projectile} is a body that is given an initial velocity and then moves under the effect
of gravity alone, like the ball in Figure~\ref{fig:projectile}. In this model we neglect air
resistance and treat $g$ as constant; we also ignore the curvature and the rotation of the Earth,
which is fine for heights and distances that are small compared with the size of the Earth. Since
the acceleration is vertical and nothing accelerates the body sideways, the motion stays in the
vertical plane that contains the initial velocity: a three-dimensional problem has become a
two-dimensional one, and we put the $x$ and $y$ axes in this plane. Let us orient the axes as in
Figure~\ref{fig:projectile} ($x$ horizontal, positive in the direction of the motion, and $y$
vertically upward), and consider a body that starts its motion with initial speed $v_0$ at an
angle $\theta$ above the horizontal. The components of its initial velocity along $x$ and $y$ are
\begin{equation}\label{eq:v0-components}
  v_{0x} = v_0\cos\theta , \qquad v_{0y} = v_0\sin\theta .
\end{equation}
Let us list what we need to know to find the position of the body in the $xy$ plane at any time
$t$. As \eqref{eq:const-a3} shows, we need the position and the velocity along $x$ and along $y$
at $t = 0$, and the (constant) acceleration along each axis. The initial velocities are given
by \eqref{eq:v0-components}. Nothing makes the body accelerate along the $x$-axis, so $a_x = 0$;
but along the $y$-axis there is the acceleration due to gravity. So the acceleration of the body
along $y$ is
\begin{equation}\label{eq:ay}
  a_y = -g .
\end{equation}
The minus sign comes from the positive direction we chose for the $y$-axis: the acceleration points
down, and we chose up as positive ($g$ itself is the positive number \qty{9.8}{m/s^2}). If we
choose the starting point as $x_0 = 0$ and $y_0 = 0$, the position and the velocity of the body
along these axes at any time $t$ are given in Table~\ref{tab:projectile}.

\begin{table}[ht]
\centering
\renewcommand{\arraystretch}{1.25}
\begin{tabular}{lcc}
\toprule
 & $x$-axis & $y$-axis \\
\midrule
At $t = 0$ & $x_0 = 0$ & $y_0 = 0$ \\
 & $v_{0x} = v_0\cos\theta$ & $v_{0y} = v_0\sin\theta$ \\
\midrule
At all times & $a_x = 0$ & $a_y = -g$ \\
\midrule
At any time $t$ & $x = (v_0\cos\theta)\,t$ & $y = (v_0\sin\theta)\,t - \frac12 gt^2$ \\
 & $v_x = v_0\cos\theta$ & $v_y = v_0\sin\theta - gt$ \\
\bottomrule
\end{tabular}
\caption{Position and velocity of a projectile at any time $t$ (launched from the origin with speed
$v_0$ at angle $\theta$; the acceleration is constant).}\label{tab:projectile}
\end{table}

As this table shows, during projectile motion the velocity along $x$ never changes. The reason, as
we said, is that the acceleration has only a $y$ component, and the motion along $y$ is completely
independent of the motion along $x$. So if you throw a ball with speed $v_0$ at an angle $\theta$
and run along the $x$ direction yourself with speed $v_0\cos\theta$, you stay exactly under the
ball and you are sure to catch it (as long as air resistance is negligible). Seen from you, the
ball simply goes straight up and comes straight back down: a first glimpse of relative motion,
which we study in Section~\ref{sec:relative}.

What is the path, or \emph{trajectory}, of this body, which starts at $x_0 = 0$, $y_0 = 0$ with
speed $v_0$ at angle $\theta$ and moves with the gravitational acceleration in the $-\jhat$
direction? To find the trajectory, that is $y$ as a function of $x$, we eliminate $t$. First, from
$x = (v_0\cos\theta)t$ we find
\begin{equation}\label{eq:t-of-x}
  t = \frac{x}{v_0\cos\theta}
\end{equation}
(we assume $\theta \neq 90^\circ$; otherwise the motion is purely vertical). Using this value of
$t$ in the equation $y = (v_0\sin\theta)t - \frac12 gt^2$ of the table, we find
\[
  y = v_0\sin\theta\,\frac{x}{v_0\cos\theta} - \frac12 g\,\frac{x^2}{(v_0\cos\theta)^2} ,
\]
that is,
\begin{equation}\label{eq:trajectory}
  y = (\tan\theta)\,x - \frac12\,\frac{g}{(v_0\cos\theta)^2}\,x^2 .
\end{equation}
This equation has the form $y = bx - cx^2$ with constants $b$ and $c > 0$: as expected, it is the
downward-opening parabola of Figure~\ref{fig:projectile}.

Now let us go a bit further and ask ourselves how high the body can get. The easiest way to answer
this question is to find when the $y$ component of the velocity is zero: if $v_y$ is not zero but
positive, the body is still moving upward, and its present height is not yet the maximum height.
Using $v_y = v_0\sin\theta - gt$, we find the time $t\sub{top}$ at which the body reaches its
highest point:
\begin{equation}\label{eq:t-top}
  0 = v_0\sin\theta - gt , \qquad \Rightarrow \qquad t\sub{top} = \frac{v_0\sin\theta}{g} .
\end{equation}
Putting this value of $t$ into $y = (v_0\sin\theta)t - \frac12 gt^2$, we find
\[
  y\sub{max} = (v_0\sin\theta)\,\frac{v_0\sin\theta}{g} - \frac12 g\Bigl(\frac{v_0\sin\theta}{g}\Bigr)^2
            = \frac{(v_0\sin\theta)^2}{g} - \frac{(v_0\sin\theta)^2}{2g} ,
\]
so the maximum height is
\begin{equation}\label{eq:ymax}
  y\sub{max} = \frac{(v_0\sin\theta)^2}{2g} .
\end{equation}
This result agrees with our expectations: if we increase the initial speed, or throw the body
upward at a steeper angle, we expect it to go higher. Only the vertical component
$v_{0y} = v_0\sin\theta$ appears: the body rises exactly as high as a ball thrown straight up with
speed $v_0\sin\theta$. Notice also that at the highest point the velocity is \emph{not} zero: its
$y$ component is zero, but its $x$ component is still $v_0\cos\theta$, so the body is moving
horizontally. And the acceleration there is still $-g\jhat$, as it is everywhere along the path.

When does the body land? There are two ways to find the answer. First, setting $y = 0$ in
$y = (v_0\sin\theta)t - \frac12 gt^2$, we find the time at which the body is back on the ground:
\begin{equation}\label{eq:t-flight}
  0 = (v_0\sin\theta)\,t - \tfrac12 gt^2 = t\bigl(v_0\sin\theta - \tfrac12 gt\bigr) ,
  \qquad \Rightarrow \qquad t = 0 \quad\text{or}\quad t\sub{tot} = \frac{2v_0\sin\theta}{g} .
\end{equation}
The first solution tells us that the body is at $y = 0$ at $t = 0$, where it started; the second,
that it is again at $y = 0$ at $t = 2v_0\sin\theta/g$, so this is when it lands. An easier way to
find the landing time is to use the symmetry of the parabola: the time it takes the body to rise
from the ground to the highest point must equal the time it takes to come down from the highest
point to the ground. So, using \eqref{eq:t-top}, we find $t\sub{tot} = 2t\sub{top} =
2v_0\sin\theta/g$. (This symmetry argument, like \eqref{eq:t-flight}, holds only if the body lands
at the same height from which it was launched. For any other landing height we must solve the
quadratic equation $y(t) = y\sub{landing}$.)

Finally, let us see where on the $x$-axis the body lands. Using $x = (v_0\cos\theta)t$ and the
landing time $t\sub{tot} = 2v_0\sin\theta/g$, we find the horizontal distance the body travels
before it lands, its \textbf{horizontal range} $R$:
\begin{equation}\label{eq:range}
  R = x\sub{max} = \frac{2v_0^2\cos\theta\sin\theta}{g} = \frac{v_0^2\sin 2\theta}{g} ,
\end{equation}
where we used the identity $2\sin\theta\cos\theta = \sin 2\theta$. Like \eqref{eq:t-flight}, this
formula holds only if the body lands at the height from which it was launched, for example a ball
kicked along level ground. For a stone thrown from a balcony, or a ball that lands on a roof, it
does not apply. Since $\sin 2\theta$ is
largest, equal to $1$, when $2\theta = 90^\circ$, the range is largest for $\theta = 45^\circ$,
where $R = v_0^2/g$. And since $\sin 2\theta = \sin(180^\circ - 2\theta)$, the two angles $\theta$
and $90^\circ - \theta$ give the same range. In real life air resistance makes the range
considerably shorter, especially for fast, light balls, and the best launch angle is then smaller
than $45^\circ$.

How should we interpret these results? Why are both the maximum height and the range proportional
to $v_0^2$? A simple argument goes like this. For a given launch angle $\theta$, the total time the
body spends in the air is proportional to its initial speed: from \eqref{eq:t-flight},
$t\sub{tot} = 2v_0\sin\theta/g \propto v_0$. On the other hand, at any given stage
of the flight (say, half way through it) the $x$ and $y$ components of the velocity are also
proportional to the initial speed. Since a distance is a velocity times a time, it is natural that
both the maximum height and the range are proportional to $v_0^2$. Dimensional analysis (Week~1)
gives the same answer even faster: the only quantities in the problem are $v_0$, $g$ and the
dimensionless angle $\theta$, and the only length we can build from $v_0$ and $g$ is $v_0^2/g$. So
every length in the problem must be $v_0^2/g$ times some function of $\theta$, exactly as in
\eqref{eq:ymax} and \eqref{eq:range}. To test these formulas with various numbers, I recommend the
\href{https://phet.colorado.edu/en/simulations/projectile-motion}{\emph{Projectile
Motion}}\footnote{\url{https://phet.colorado.edu/en/simulations/projectile-motion}} simulation of
PhET, developed at the University of Colorado.

\begin{example}[Same range, two angles]
A ball is launched from level ground at $v_0 = \qty{20}{m/s}$, once at $\theta = 30^\circ$ and once
at $\theta = 60^\circ$. From \eqref{eq:range}, both shots have the same range,
$R = (\qty{20}{m/s})^2\sin 60^\circ/(\qty{9.8}{m/s^2}) = \qty{35}{m}$. From \eqref{eq:ymax} the
maximum heights are $(\qty{10}{m/s})^2/(2\times\qty{9.8}{m/s^2}) = \qty{5.1}{m}$ and
$(\qty{17.3}{m/s})^2/(2\times\qty{9.8}{m/s^2}) = \qty{15}{m}$, and from \eqref{eq:t-flight} the
times of flight are \qty{2.0}{s} and \qty{3.5}{s}. The steeper shot goes three times higher
($\sin^2 60^\circ/\sin^2 30^\circ = 3$) and stays longer in the air, but lands at the same place.
\end{example}

\begin{figure}[ht]
\centering
\begin{tikzpicture}[>=Stealth,scale=0.95]
  % ground and tree
  \fill[pattern=north east lines,pattern color=fzgray] (-0.4,-0.3) rectangle (10.6,0);
  \draw[thick] (-0.4,0) -- (10.6,0);
  \draw[line width=3pt,fzgray] (8.1,0) -- (8.1,6.0);
  \draw[line width=1.5pt,fzgray] (8.1,5.35) -- (7.3,5.52);
  % zookeeper (stick figure) with dart gun
  \draw[thick,fzgray] (0.15,1.42) circle (0.12);
  \draw[thick,fzgray] (0.15,1.30) -- (0.15,0.55) (0.0,0) -- (0.15,0.55) -- (0.3,0)
    (0.15,1.12) -- (0.42,0.96);
  \draw[line width=2pt] (0.2,0.84) -- (0.5,1);
  \node[below,font=\small] at (0.2,-0.3) {zookeeper};
  % reference line at gun height
  \draw[fzgray] (0.5,1) -- (10.0,1);
  % line of aim and the dart's parabola
  \draw[dashed,thick] (0.5,1) -- (7.7,5);
  \draw[dashed,thick,fzcomp,domain=0:7.2,samples=40] plot ({0.5+\x},{1+0.55556*\x-0.034722*\x*\x});
  \draw[->,thick,fzcomp] (5.4,2.910) -- (5.6,2.952);
  % initial velocity
  \draw[dashed,fzgray] (1.898,1) -- (1.898,1.777) -- (0.5,1.777);
  \draw[->,very thick,fzvec] (0.5,1) -- (1.898,1.777) node[above]{$\vec{v}_0$};
  \draw[->,thick,fzcomp] (0.5,1) -- (1.898,1) node[below]{$v_{0x}$};
  \draw[->,thick,fzcomp] (0.5,1) -- (0.5,1.777) node[left]{$v_{0y}$};
  % monkey: at t = 0 (hanging from the branch) and at time t
  \draw[thick,fzaux] (7.7,5.17) -- (7.7,5.43);
  \draw[thick,fzaux,fill=white] (7.7,5) circle (0.17);
  \fill[fzaux] (7.7,3.2) circle (0.17);
  \draw[->,fzaux] (7.7,4.75) -- (7.7,3.45);
  \node[anchor=east,font=\small,fzaux] at (7.4,5.3) {monkey at $t=0$};
  \node[anchor=north east,font=\small,fzaux] at (7.55,2.75) {monkey at $t$};
  % dimensions
  \draw[fzgray] (7.9,5) -- (10.2,5);
  \draw[fzgray] (7.9,3.2) -- (9.1,3.2);
  \draw[<->] (8.9,5) -- (8.9,3.2) node[midway,right]{$\frac12 gt^2$};
  \draw[<->] (8.9,3.2) -- (8.9,1) node[midway,right]{$y$};
  \draw[<->] (10.0,5) -- (10.0,1) node[midway,right]{$h = v_{0y}t$};
  \draw[<->] (0.5,0.45) -- (7.7,0.45) node[midway,fill=white,inner sep=1pt]{$x$};
\end{tikzpicture}
\caption{The hunter and the monkey. In a time $t$ the dart reaches the tree, a horizontal distance
$x$ away. Without gravity it would be at height $h = v_{0y}t$ on the line of aim; with gravity it is
$\frac12 gt^2$ lower, at $y = v_{0y}t - \frac12 gt^2 = h - \frac12 gt^2$. In the same time the
monkey, released at $t = 0$, has fallen exactly $\frac12 gt^2$.}\label{fig:monkey}
\end{figure}

\textbf{The hunter and the monkey.} Our last example on this topic is about a monkey that has
escaped from the zoo and climbed a tree, and a zookeeper who tries to hit it with a tranquilizer
dart. The zookeeper aims directly at the monkey. If the monkey stayed where it is, the dart would
miss it: as we have just seen, gravity bends the dart's path below the straight line of aim. But
the monkey has not taken FIZ~101. The moment it sees the dart leave the gun, it lets go of the
branch and starts to fall. Is the monkey safe now? Or will the zookeeper put the poor monkey to
sleep and catch it?

After the zookeeper aims at the monkey and fires, at any time $t$ the dart would be on the straight
line of aim if there were no gravity, while with gravity it is on the parabola below it
(Figure~\ref{fig:monkey}). But here we can use an important fact that we learned when we studied
projectile motion: the motion of the dart along $x$ is not affected by gravity. So, with or without
gravity, at a given time the dart is at the same $x$. Gravity affects only the dart's $y$ position.
By how much does gravity change it? To find out, we look at
\begin{equation}\label{eq:dy-def}
  \Delta y = y_{g=0} - y ,
\end{equation}
the difference between the height the dart would have without gravity and its actual height. From
Table~\ref{tab:projectile}, $y_{g=0} = (v_0\sin\theta)t$ and $y = (v_0\sin\theta)t - \frac12 gt^2$,
so this difference is
\begin{equation}\label{eq:dy}
  \Delta y = \frac12 gt^2 .
\end{equation}
So when the dart reaches the vertical line along which the monkey falls, the monkey, to escape,
would have to have fallen more or less than $\frac12 gt^2$ below its starting point. But the monkey
let itself fall from rest at the very instant the dart was fired, so in the same time $t$ it has
fallen exactly $\frac12 gt^2$ (the $y$ equation of \eqref{eq:const-a3} with $v_{0y} = 0$ and
$a_y = -g$ gives $y - y_0 = -\frac12 gt^2$): exactly as far as
the dart has dropped below the line of aim. The dart and the monkey arrive at the same point at the
same time. The poor monkey, not knowing FIZ~101, becomes the victim of the dart.

Notice that the answer depends neither on the speed of the dart nor on the value of $g$: it would
be the same on the Moon. The only condition is that the dart be fast enough to reach the tree
before the monkey reaches the ground. There is an even quicker way to see the result. Seen from the
falling monkey, the monkey and the dart have the same downward acceleration $g$, so relative to the
monkey the dart has no acceleration at all: it flies along the straight line along which it was
aimed, straight at the monkey.

%----------------------------------------------------------------------
\subsection{Circular motion}\label{sec:circular}

As we saw in Section~\ref{sec:invariants}, rotating the coordinate system does not change the length
of a vector: the length is invariant. This is a statement about describing the \emph{same} vector
with different axes. Something different happens when the vector itself turns: if the velocity
vector of a body turns through some angle while keeping its length, the speed, which is the
magnitude of the velocity, does not change. But, as we stressed at the start of this chapter, the
velocity itself \emph{has} changed, since its direction is different; and a change of velocity is
possible only when there is an acceleration. In this section our aim is to get an idea of how to
study such motions and to calculate the acceleration in them.

\subsubsection{Uniform circular motion}\label{sec:ucm}

\begin{figure}[ht]
\centering
\begin{minipage}[c]{0.50\textwidth}
\centering
\begin{tikzpicture}[>=Stealth,scale=0.95]
  \draw[thick] (0,0) circle (2);
  \draw[->] (-2.4,0) -- (2.7,0) node[right]{$x$};
  \draw[->] (0,-2.4) -- (0,2.7) node[above]{$y$};
  % position vector at 35 degrees: P = (1.638,1.147)
  \draw[->,very thick] (0,0) -- (1.638,1.147);
  \node[right] at (1.9,0.95) {$\vec{r} = r\hat{r}$};
  \draw (0.6,0) arc[start angle=0,end angle=35,radius=0.6];
  \node at (0.95,0.24) {$\omega t$};
  % velocity, tangent
  \draw[->,very thick,fzvec] (1.638,1.147) -- (0.892,2.212) node[above left=-2pt]{$\vec{v}$};
  % acceleration, towards the center (drawn just beside r, on its upper side)
  \draw[->,very thick,fzaux] (1.523,1.311) -- (0.622,0.680);
  \draw[fzaux] (0.95,0.93) -- (-0.9,2.45) node[left,font=\small]{$\vec{a} = -\omega^2 r\,\hat{r}$};
  \fill (1.638,1.147) circle (1.8pt);
\end{tikzpicture}
\end{minipage}\hfill
\begin{minipage}[c]{0.42\textwidth}
\caption{Uniform circular motion. At time $t$ the position vector makes an angle $\omega t$ with
the $x$-axis. The velocity is tangent to the circle, and the acceleration points towards the
center.}\label{fig:ucm}
\end{minipage}
\end{figure}

The motion of a body around a circle with constant speed is called \textbf{uniform circular
motion}. Before we start studying it, it is useful to make some definitions.
\begin{itemize}
  \item \textbf{Period ($T$):} the time it takes the body to go once around the circle.
  \item \textbf{Frequency ($f$):} the number of revolutions the body makes per second. The unit of
    frequency is $\unit{s^{-1}}$, more commonly called the hertz (\unit{Hz}). The frequency is
    \begin{equation}\label{eq:f}
      f = \frac{1}{T} .
    \end{equation}
  \item \textbf{Angular velocity ($\omega$):} the angle, in radians, that the body turns through
    per second; its unit is \unit{rad/s}. Since one full turn is $2\pi$ radians, the angular
    velocity is
    \begin{equation}\label{eq:omega}
      \boxed{\,\omega = \frac{2\pi}{T}\,} \;=\; 2\pi f .
    \end{equation}
    (Here $\omega$ is a positive number, strictly the angular \emph{speed}; when we study rotation
    in Chapter~10 we will see that angular velocity also has a direction.)
\end{itemize}
With these definitions, if the body moves around a circle of radius $r$ as in Figure~\ref{fig:ucm},
its speed $v$ is
\begin{equation}\label{eq:v-omega-r}
  \boxed{\,v = \frac{2\pi r}{T} = \omega r\,} ,
\end{equation}
since in one period it travels once around the circumference, $2\pi r$. Although the speed of the
body does not change during the motion, its velocity does, so there must be an acceleration. Let
the angular velocity of the body be $\omega$, and let us choose the axes so that the body is on the
positive $x$-axis (on the line $y = 0$) at $t = 0$ and moves counterclockwise. Then at time $t$ its
position vector makes an angle $\omega t$ with the $x$-axis, as in Figure~\ref{fig:ucm}, and the
position of the body is
\begin{equation}\label{eq:ucm-r}
  \vec{r} = \ihat\, r\cos(\omega t) + \jhat\, r\sin(\omega t) .
\end{equation}
It is convenient to write this as $\vec{r} = r\hat{r}$, where
$\hat{r} = \ihat\cos(\omega t) + \jhat\sin(\omega t)$ is the unit vector that points from the center
towards the body. Unlike $\ihat$ and $\jhat$, the unit vector $\hat{r}$ turns as the body moves.
We differentiate \eqref{eq:ucm-r} with respect to time, with $r$ and $\omega$ constant, using the
chain rule: $\dd\cos(\omega t)/\dd t = -\omega\sin(\omega t)$ and
$\dd\sin(\omega t)/\dd t = \omega\cos(\omega t)$. The velocity of the body is then
\begin{equation}\label{eq:ucm-v}
  \vec{v} = -\ihat\, \omega r\sin(\omega t) + \jhat\, \omega r\cos(\omega t) ,
\end{equation}
and, differentiating once more, its acceleration is
\begin{equation}\label{eq:ucm-a}
  \vec{a} = -\ihat\, \omega^2 r\cos(\omega t) - \jhat\, \omega^2 r\sin(\omega t) = -\hat{r}\,\omega^2 r
          = -\omega^2\vec{r} .
\end{equation}
From \eqref{eq:ucm-v}, the speed is $v = |\vec{v}| = \omega r\sqrt{\sin^2(\omega t) +
\cos^2(\omega t)} = \omega r$, as we expected, and $\vec{v}\cdot\vec{r} = 0$: the velocity is
perpendicular to the radius, that is, tangent to the circle. The acceleration \eqref{eq:ucm-a}
points along $-\hat{r}$, from the body towards the center of the circle; it is called
\textbf{centripetal} (``center-seeking'') acceleration. Its magnitude is
\begin{equation}\label{eq:arad}
  a\sub{rad} = \omega^2 r = \frac{v^2}{r} ,
\end{equation}
where we used $\omega = v/r$ from \eqref{eq:v-omega-r}. We use the subscript ``rad'' because, as
\eqref{eq:ucm-a} shows, the acceleration at every point is along the radius, in the radial
direction, pointing inward. Using $v = 2\pi r/T$, we can also write it as
$a\sub{rad} = 4\pi^2 r/T^2$.

Two warnings. First, the magnitude $a\sub{rad}$ is constant, but the acceleration \emph{vector} is
not: its direction turns with the body. So the constant-acceleration equations \eqref{eq:const-a3}
do not apply to circular motion. Second, the acceleration is not along the velocity: in uniform
circular motion it is perpendicular to it, and this is exactly why it changes only the direction of
the velocity and not the speed. Kinematics only tells us that this inward acceleration must be
present; something must provide it: the tension in a string for a stone whirled around, friction
on the tires for a car taking a curve, gravity for the Moon. When we study Newton's laws we will
see that a body in uniform circular motion needs a net force of magnitude $mv^2/r$ pointing towards
the center.

\begin{example}[Standing on the equator]
Because the Earth spins, a person standing on the equator moves around a circle of radius
$r = \qty{6.37e6}{m}$ once every $T = \qty{23}{h}\,\qty{56}{min} = \qty{8.62e4}{s}$ (one rotation
of the Earth with respect to the stars). The person's speed is
\[
  v = \frac{2\pi r}{T} = \frac{2\pi(\qty{6.37e6}{m})}{\qty{8.62e4}{s}} = \qty{464}{m/s} ,
\]
about \qty{1670}{km/h}, and the centripetal acceleration is $a\sub{rad} = v^2/r = \qty{0.034}{m/s^2}$,
only about \qty{0.35}{\percent} of $g$. This is one of the reasons why the measured value of $g$ is
slightly smaller at the equator than at the poles.
\end{example}

\subsubsection{Nonuniform circular motion}\label{sec:nucm}

In the previous section we discussed circular motion in which the speed of the body does not
change. If the speed of the body changes during the motion, as for a car speeding up around a curve
or a stone whirled on a string in a vertical circle, we call it \textbf{nonuniform circular
motion}. In nonuniform circular motion the acceleration has both a radial component and a tangential
component. Since the radial direction and the direction tangent to the circle are perpendicular to
each other, any acceleration in the plane of the circle can be written in terms of vectors along
these two directions (Figure~\ref{fig:nucm}).

\begin{figure}[ht]
\centering
\begin{minipage}[c]{0.46\textwidth}
\centering
\begin{tikzpicture}[>=Stealth,scale=1.0]
  \draw[thick] (0,0) circle (2);
  \fill (0,0) circle (1.2pt);
  % point at 35 degrees: P = (1.638,1.147)
  \draw[fzgray] (0,0) -- (1.638,1.147);
  % velocity, along the counterclockwise tangent (drawn first, under a_tan)
  \draw[->,thick,fzvec] (1.638,1.147) -- (0.720,2.457) node[above left=-1pt]{$\vec{v}$};
  % a_rad: towards the center, length 1.1
  \draw[->,very thick,fzaux] (1.638,1.147) -- (0.737,0.516) node[below right=-1pt]{$\vec{a}\sub{rad}$};
  % a_tan: along v (speeding up), length 0.8
  \draw[->,very thick,fzaux] (1.638,1.147) -- (1.179,1.802);
  \node[fzaux,anchor=west] at (1.5,1.68) {$\vec{a}\sub{tan}$};
  % total
  \draw[dashed,fzgray] (0.737,0.516) -- (0.278,1.171) -- (1.179,1.802);
  \draw[->,very thick,fzcomp] (1.638,1.147) -- (0.278,1.171) node[left]{$\vec{a}$};
  % direction of motion
  \draw[->,thick,fzvec] (2.25,-0.55) arc[start angle=-14,end angle=14,radius=2.3];
  \fill (1.638,1.147) circle (1.8pt);
\end{tikzpicture}
\end{minipage}\hfill
\begin{minipage}[c]{0.46\textwidth}
\caption{Nonuniform circular motion of a body that is speeding up (moving counterclockwise). The
radial component of the acceleration points towards the center, the tangential component points
along the velocity $\vec{v}$, and their sum $\vec{a}$ no longer points towards the
center.}\label{fig:nucm}
\end{minipage}
\end{figure}

To see where the two components come from, let $\phi(t)$ be the angle the position vector makes
with the $x$-axis, so that $\vec{r} = r(\ihat\cos\phi + \jhat\sin\phi)$ with $r$ constant, and let
$\omega = \dd\phi/\dd t$, which now changes with time (we take the body to move counterclockwise,
so $\omega > 0$). By the chain rule, $\vec{v} = \omega r(-\ihat\sin\phi + \jhat\cos\phi)$: the
velocity is still tangent to the circle, with speed $v = \omega r$. Differentiating once more, now
with the product rule because $\omega$ depends on $t$,
\[
  \vec{a} = r\dv{\omega}{t}\,(-\ihat\sin\phi + \jhat\cos\phi) - \omega^2 r\,(\ihat\cos\phi + \jhat\sin\phi) .
\]
The second term points towards the center and is exactly the centripetal acceleration of
\eqref{eq:ucm-a}. The first points along the tangent, and its component
$r\,\dd\omega/\dd t = \dd v/\dd t$ is the rate of change of the speed. (For uniform motion
$\dd\omega/\dd t = 0$ and $\phi = \omega t$, and we are back to \eqref{eq:ucm-a}.)

The radial component, as we discussed in the previous section, changes the direction of the
velocity; it points towards the center and its magnitude is
\begin{equation}\label{eq:arad2}
  a\sub{rad} = \omega^2 r = \frac{v^2}{r} ,
\end{equation}
where now $\omega$ and $v$ are the values at that instant, which change with time. The tangential
component changes the speed of the body, and is given by
\begin{equation}\label{eq:atan}
  a\sub{tan} = \frac{\dd|\vec{v}|}{\dd t} .
\end{equation}
We use the subscript ``tan'' to stress that this component of the acceleration is tangent to the
circle, that is, along the line of the velocity. It is positive (in the direction of $\vec{v}$) when the body
speeds up, and negative (opposite to $\vec{v}$) when it slows down. The total acceleration is the
vector sum of the two components, $\vec{a} = \vec{a}\sub{rad} + \vec{a}\sub{tan}$, with magnitude
$a = \sqrt{a\sub{rad}^2 + a\sub{tan}^2}$; unlike in uniform circular motion, it no longer points
towards the center. (In uniform circular motion $a\sub{tan} = 0$, and we are back to
Section~\ref{sec:ucm}.)

%----------------------------------------------------------------------
\subsection{Relative velocity}\label{sec:relative}

In principle, every observer who has a meter stick and a stopwatch, that is, who can make
measurements, defines a \textbf{frame of reference}: a frame of reference consists of a coordinate
system (an origin and axes, attached to some body such as the ground, a car or a train) together
with a clock that measures time.

The values of the physical quantities we studied in the last few sections (position, velocity and
acceleration) depend on the frame in which we give them. For example, if you are traveling in a car
with a friend, you measure the velocity of your friend sitting next to you as $\vec{0}$, while an
observer outside the car, waiting at a bus stop for example, measures your friend's velocity to be
the velocity of the car. It is therefore important to know how different frames agree with each
other, that is, how a measurement made in one frame is converted into another frame. Throughout, we
assume, as Newton did, that all observers use the same time: a clock in the car and a clock at the
bus stop tick at the same rate. This is true to excellent accuracy when all speeds are small
compared with the speed of light; near the speed of light we would need Einstein's special theory
of relativity.

\begin{figure}[ht]
\centering
\begin{minipage}[c]{0.52\textwidth}
\centering
\begin{tikzpicture}[>=Stealth,scale=0.95]
  % frame A
  \draw[thick] (0,-0.3) -- (0,2.4) node[above]{$y_A$};
  \draw[->,thick] (-0.4,0) -- (5.8,0) node[right]{$x_A$};
  \node[below left] at (0,0) {$O_A$};
  % frame B
  \draw[thick,fzaux] (2.4,-0.3) -- (2.4,2.4) node[above]{$y_B$};
  \draw[very thick,fzaux] (2.0,0) -- (5.5,0);
  \node[above left,fzaux] at (2.4,0) {$O_B$};
  \draw[->,very thick,fzvec] (2.4,1.6) -- (3.6,1.6) node[right]{$\vec{v}_{B/A}$};
  % point P
  \fill (4.9,0) circle (2pt) node[above]{$P$};
  % dimensions
  \draw[|<->|] (0,-0.75) -- (2.4,-0.75) node[midway,fill=white,inner sep=1pt]{$x_{B/A}$};
  \draw[|<->|,fzaux] (2.4,-0.75) -- (4.9,-0.75) node[midway,fill=white,inner sep=1pt]{$x_{P/B}$};
  \draw[|<->|] (0,-1.35) -- (4.9,-1.35) node[midway,fill=white,inner sep=1pt]{$x_{P/A}$};
\end{tikzpicture}
\end{minipage}\hfill
\begin{minipage}[c]{0.42\textwidth}
\caption{Frame $B$ moves along the common $x$-axis relative to frame $A$ with velocity
$\vec{v}_{B/A}$. The position of the point $P$ is measured in both frames.}\label{fig:relative}
\end{minipage}
\end{figure}

Let us first study the simpler problem in one dimension. As in Figure~\ref{fig:relative}, let
$x_{P/A}$ be the position of a point $P$ in frame $A$ (read: ``the position of $P$ relative to
$A$''), and let $x_{P/B}$ be its position in a frame $B$ whose $x$-axis lies along the same line and
points in the same direction. The position of the origin of frame $B$ relative to frame $A$ is
$x_{B/A}$. Then, as the figure shows, the position $x_{P/A}$ is
\begin{equation}\label{eq:x-rel}
  x_{P/A} = x_{P/B} + x_{B/A} .
\end{equation}
If we write $\relx{v}{P/A}$ for the velocity of $P$ relative to frame $A$ (and similarly for the
other subscripts), we find the relation between the velocities by taking the time derivative of
\eqref{eq:x-rel}:
\begin{equation}\label{eq:v-rel}
  \dv{x_{P/A}}{t} = \dv{x_{P/B}}{t} + \dv{x_{B/A}}{t} , \qquad \Rightarrow \qquad
  \relx{v}{P/A} = \relx{v}{P/B} + \relx{v}{B/A} .
\end{equation}
(A handy way to remember it: the ``inner'' subscripts $B$ cancel, so that $P/B$ followed by $B/A$
gives $P/A$.) A simple consequence of these equations is the following. Suppose the point $P$ is
the origin $O_A$ of frame $A$ itself, so that the observer in $B$ is looking at the origin of $A$.
Then $x_{P/A} = 0$ (the origin of $A$ is at zero in its own frame) and $x_{P/B} = x_{A/B}$ (the
position of $O_A$ measured in $B$), and \eqref{eq:x-rel} gives $0 = x_{A/B} + x_{B/A}$, that is,
\begin{equation}\label{eq:x-AB}
  x_{A/B} = -x_{B/A} .
\end{equation}
Taking the derivative of this equation, we find
\begin{equation}\label{eq:v-AB}
  \relx{v}{A/B} = -\relx{v}{B/A} .
\end{equation}
This result agrees with what we expect. If, waiting at a bus stop, we measure the velocity of a
vehicle as $v$ in the $+x$ direction, an observer inside the vehicle sees us, waiting at the stop,
move with velocity $-v$, in the other direction.

\begin{example}[Overtaking and oncoming cars]
You ($B$) drive along a straight road at $\relx{v}{B/E} = \qty{+25}{m/s}$ relative to the Earth
($E$). A car $P$ ahead of you moves at $\relx{v}{P/E} = \qty{+30}{m/s}$. From \eqref{eq:v-rel},
$\relx{v}{P/B} = \relx{v}{P/E} - \relx{v}{B/E} = \qty{+5}{m/s}$: you see it pull slowly away. An
oncoming car with $\relx{v}{P/E} = \qty{-30}{m/s}$ has $\relx{v}{P/B} = -30 - 25 = \qty{-55}{m/s}$
relative to you, which is why oncoming traffic seems to rush past. And by \eqref{eq:v-AB}, the road
moves backward relative to you: $\relx{v}{E/B} = \qty{-25}{m/s}$.
\end{example}

We can generalize the idea of relative velocity to motion in three dimensions, just as we
generalized straight-line motion to three dimensions using vectors. Now we speak of the position
vector $\vec{r}_{P/A}$ of the point $P$ relative to frame $A$. We take the axes of $A$ and $B$
parallel to each other, so that $B$ moves relative to $A$ without rotating. Generalizing
\eqref{eq:x-rel} to three dimensions, the position $\vec{r}_{P/A}$ is
\begin{equation}\label{eq:r-rel}
  \vec{r}_{P/A} = \vec{r}_{P/B} + \vec{r}_{B/A} .
\end{equation}
Similarly, by differentiating, we find the relation between the velocities:
\begin{equation}\label{eq:galileo}
  \vec{v}_{P/A} = \vec{v}_{P/B} + \vec{v}_{B/A} .
\end{equation}
This transformation is called the \textbf{Galilean velocity transformation}. As in straight-line
motion, if we take $P$ to be the origin of frame $A$, we find, as expected,
\begin{equation}\label{eq:v-AB3}
  \vec{v}_{A/B} = -\vec{v}_{B/A} .
\end{equation}
If $B$ moves relative to $A$ with constant velocity, $\vec{v}_{B/A}$ does not change, and
differentiating \eqref{eq:galileo} once more gives $\vec{a}_{P/A} = \vec{a}_{P/B}$: observers in
the two frames measure \emph{different} velocities for $P$, but the \emph{same} acceleration. This
simple fact will be the key to Newton's laws and to inertial frames of reference in Chapter~5,
where we come back to Galilean relativity.

\begin{example}[Crossing a river]
A boat ($B$) heads straight across a river, in the $+y$ direction, at \qty{4.0}{m/s} relative to the
water ($W$), and the river flows in the $+x$ direction at \qty{3.0}{m/s} relative to the Earth
($E$). From \eqref{eq:galileo}, the velocity of the boat relative to the Earth is
\[
  \vec{v}_{B/E} = \vec{v}_{B/W} + \vec{v}_{W/E} = (\qty{3.0}{m/s})\,\ihat + (\qty{4.0}{m/s})\,\jhat ,
\]
with magnitude $\sqrt{3.0^2 + 4.0^2}\,\unit{m/s} = \qty{5.0}{m/s}$, at an angle
$\arctan(3.0/4.0) = 37^\circ$ downstream from the straight-across direction. The boat still crosses
in the same time as it would in still water, because the current changes only the $x$ component of
its velocity: once again, the components are independent.
\end{example}

%======================================================================
\section*{About these notes}
These notes are not original. They draw heavily on the books listed in the Week~1 notes, on Walter Lewin's 8.01x Classical Mechanics lectures at MIT, and on the Dynamics and Relativity courses taught at different times by Stephen Siklos and David Tong at the University of Cambridge. They are not official course notes of FIZ~101: they are written only to help the students of Mehmet \"Ozkan's section. Istanbul Technical University and its Department of Physics Engineering bear no responsibility for them.

\end{document}
